\subsection{The chart family with two-edge circles}

\begin{lemma}[Boundary values distinguish degenerate cells]
\label{C:retraction-boundary-injective}
For the commuting retractions in
Lemma~\ref{C:two-retraction-cubical-model}, the map
\[
 U\cup V\longrightarrow T^4,\qquad
 t\longmapsto(P_0t,P_1t,Q_0t,Q_1t)
\]
is injective.
\end{lemma}
\begin{proof}
Two elements of $U$ are distinguished by $P_0$, and two
elements of $V$ by $Q_0$.  If $x\in U$, $y\in V$ have the
same boundary tuple, then $P_0y=x$ and $Q_0x=y$.
Commutation gives
\[
 x=P_0Q_0x=Q_0P_0x=Q_0x=y.
\]
These cases exhaust the union.
\end{proof}

\begin{theorem}[The two-edge-circle chart family has zero products]
\label{C:affine-four-circle-family-zero}
Keep the corner monoids and support of
Theorem~\ref{C:affine-chart-family-zero}, but take the two
adjacent hom-sets to be
\[
 A=\{u_0,u_1,a_0,a_1\},\qquad
 B=\{v_0,v_1,b_0,b_1\}.
\]
The two nonidentity precomposition maps fix both $u_\eta$
and send each $a_j$ to $u_0$ or $u_1$, respectively.
The two nonidentity postcomposition maps fix both $v_\eta$
and send each $b_i$ to $v_0$ or $v_1$, respectively.
Allow any finite nonempty long-edge biset $T=C(0,2)$ and
any equivariant composition.  Every rational product of
two degree-two classes on $BC$ is zero.  Thus this entire
family fails to realize $\mathbb {CP}^2$ at every size.
\end{theorem}
\begin{proof}
The two adjacent action graphs are circles, each with two
vertices and two equally oriented parallel edges.  Write
$\Gamma_A,\Gamma_B$ for them.  As in the proof of
Theorem~\ref{C:affine-chart-family-zero}, the natural graph
and cubical models give a triangle of spaces with contractible
vertices, edges $\Gamma_A,\Gamma_B,K(T)$, and face
$\Gamma_A\times\Gamma_B$.  Its face maps are the two
projections and the cellular composition map
\[
 h:\Gamma_A\times\Gamma_B\longrightarrow K(T).
\]
Put $t_{ij}=b_i a_j$.  Equivariance gives, for every
$\varepsilon,\delta\in\{0,1\}$,
\begin{equation}\label{C:four-circle-boundary-equalities}
 P_\varepsilon t_{i0}=P_\varepsilon t_{i1},\qquad
 Q_\delta t_{0j}=Q_\delta t_{1j}.
\end{equation}
The remaining composition values are their boundary values.

Let $[t]$ denote the two-cell of $K(T)$ indexed by $t$ when
$t\notin U\cup V$, and put $[t]=0$ for a degenerate cell.
The fundamental cellular two-cycle of the product of the
two circles maps to
\begin{equation}\label{C:four-circle-top-chain}
 z=[t_{00}]-[t_{01}]-[t_{10}]+[t_{11}].
\end{equation}
This is a cycle because $h$ is a cellular map.  If $z\ne0$
as a cellular chain, it is nonzero in $H_2(K(T);\mathbb Q)$,
since the target has no three-cells.  The only total-degree-four
term in the triangle spectral sequence is
$H_2(\Gamma_A\times\Gamma_B;\mathbb Q)=\mathbb Q$.
Its first differential sends its generator to $z$; there
is no other nonzero term in total degree four.  Thus
$H_4(BC;\mathbb Q)=0$ in this case.

Suppose now that $z=0$.  First assume that one of its four
states is nondegenerate.  Cancellation in the free cellular
two-chain group forces a positive occurrence to equal a
negative occurrence as an actual state.  Such positions
are adjacent in the two-by-two array.  Interchanging edge
names, and exchanging the two coordinates in this argument
if necessary, reduces to
\[
 t_{00}=t_{01}\notin U\cup V.
\]
Equation~\eqref{C:four-circle-top-chain} gives
$[t_{10}]=[t_{11}]$.  If this common cell is nonzero, its
indices are equal.  Otherwise both states belong to $U\cup V$.
Their $P$-boundaries agree by
\eqref{C:four-circle-boundary-equalities}; their $Q$-boundaries
agree by the same equations and $t_{00}=t_{01}$.
Lemma~\ref{C:retraction-boundary-injective} again gives
$t_{10}=t_{11}$.  Thus the entire two-by-two array has equal
columns.  All its boundary values agree accordingly, so
$h$ factors through the quotient of $\Gamma_A$ identifying
its two equally oriented edges.  That quotient is an interval.
Contracting the interval to its initial vertex gives a based
homotopy of $h$ to a map factoring through
$\operatorname{pr}_{\Gamma_B}$.  The equal-row case gives
the other projection.

It remains to consider the case where all four $t_{ij}$
are degenerate.  The image of $h$ then lies in the one-skeleton
of $K(T)$, a graph.  The two circle generators of the source
map to closed edge paths in its two distinct edge colors:
one uses only the cells of $U\setminus V$, and the other
uses only the cells of $V\setminus U$.  These based loops
commute in the graph fundamental group, because they come
from a map of a torus.  They cannot both be nontrivial.
Indeed, reduce each edge path by cancellation of an edge
with its inverse.  If both reduced paths were nonempty,
the word formed by their commutator would have no cancellation
at any of its four junctions: adjacent factors have disjoint
edge colors.  It would therefore be a nontrivial reduced
loop in a graph, a contradiction.

At least one coordinate loop is consequently null-homotopic.
The connected graph component containing the image is a
$K(F,1)$.  The map $h$ is based homotopic to the map obtained
by restricting to the other coordinate circle and then
projecting onto it, since these maps induce the same
homomorphism on fundamental groups.  This conclusion also
covers the case where both loops are trivial.

We have proved that whenever $z=0$, the face map $h$ is
homotopic to a map through one of the adjacent-edge
projections.  Replacing this face map by that homotopic map
preserves the triangle homotopy colimit; the maps to its
vertex spaces are constant, so no additional compatibility
is required.  The suspension construction in the proof of
Theorem~\ref{C:affine-chart-family-zero} now gives a map
from this homotopy colimit to $\Sigma K(T)$ inducing an
isomorphism on rational degree-two cohomology.  Explicitly,
if $h=\alpha\operatorname{pr}_{\Gamma_B}$, use the identity
suspension on the long edge, $\Sigma\alpha$ on the $B$ edge,
and collapse the $A$ edge; the barycentric coordinate of
their common extreme vertex supplies the face formula.
The degree-two group is the kernel of
\[
 H^1(\Gamma_A)\oplus H^1(\Gamma_B)\oplus H^1(K(T))
 \longrightarrow H^1(\Gamma_A\times\Gamma_B),
\]
and projection of this kernel onto $H^1(K(T))$ is an
isomorphism.  The edge-to-face degree-zero map is onto,
so there is no additional degree-two term, even when
$K(T)$ is disconnected.  Since positive products on a
suspension vanish, every claimed product on $BC$ vanishes.
\end{proof}

\begin{remark}[Completed controls through twenty-one arrows]
\label{C:affine-four-circle-controls}
The exact generator \path{code/d_affine_four_circles.py}
enumerates all commuting retraction bisets and all equivariant
composition maps for $1\leq|T|\leq6$, using only actual
target-state and transient-edge automorphisms to remove
repetitions.  The respective arrow totals and numbers of
literal $H_1=0$, $\dim H_2=1$ mod-two candidates are
\[
\begin{array}{c|rrrrrr}
\text{arrows}&16&17&18&19&20&21\\\hline
\text{candidates}&0&6&44&340&2414&19652.
\end{array}
\]
The independent bar checker verifies a zero-square primitive
for all $22456$ candidates, including $4092904$ typed
associativity equations.  All six generator and checker
streams are complete; their prefixes are
\path{code/outputs/d_affine_four_circles_1} through
\path{code/outputs/d_affine_four_circles_6}.
These computations corroborate the theorem and are not used
in its all-size proof.
\end{remark}
